% !TEX program = lualatex
% ECON 201, Worksheet 16: A Firm That Takes the Price as Given.
% Compile with LuaLaTeX so the PDF is tagged. Build from the course root:
%   latexmk -lualatex -cd worksheets/worksheet16.tex && latexmk -c -cd worksheets/worksheet16.tex
% Every number here is computed and printed by slides/figures/make_lecture16.py
% (Lena's lemonade stand, its own numbers rather than the deck's; the
% script also draws figures/lemonade.pdf).
\DocumentMetadata{
  pdfstandard=UA-2,
  pdfversion=2.0,
  lang=en-US,
  tagging=on
}
\documentclass[11pt]{article}

\usepackage[letterpaper, margin=0.75in, top=0.5in, bottom=0.45in]{geometry}
\usepackage{fontspec}
\setmainfont{Fira Sans}
\usepackage{amsmath}
\usepackage{amssymb}
\usepackage{graphicx}
\usepackage{booktabs}
\usepackage{array}
\usepackage{xcolor}
\usepackage{needspace}
\usepackage[most]{tcolorbox}
\definecolor{cream}{HTML}{FDF3EA}
\definecolor{accent}{HTML}{BF5700}
\usepackage{titlesec}
\titleformat{\section}{\large\bfseries\color{accent}}{}{0pt}{}
\titlespacing*{\section}{0pt}{4pt}{2pt}
\usepackage{hyperref}
\hypersetup{pdftitle={ECON 201 Worksheet: A Firm That Takes the Price as Given}, pdfauthor={Div Bhagia}, hidelinks}
\pagestyle{empty}
\setlength{\parindent}{0pt}
\setlength{\parskip}{2pt}

% Ruled answer space: n lines. Plain TeX loop; pgffor's \foreach breaks the
% enumerate counter when used inside an item.
\newcount\answerlinecount
\newcommand{\answerlines}[1]{%
  \par\vspace{1pt}\answerlinecount=#1
  \loop\ifnum\answerlinecount>0
    \rule{\linewidth}{0.3pt}\par\vspace{5pt}%
    \advance\answerlinecount by -1
  \repeat
  \vspace{-5pt}}

% Ruled grid rather than booktabs: students write in the cells, so every
% cell gets visible borders and the blank cells get extra height.
\newcolumntype{C}{>{\centering\arraybackslash}p{3.1cm}}
\newcommand{\blankrow}{\rule{0pt}{1.6em}}

\begin{document}

{\LARGE\bfseries\color{accent} Worksheet: Taking the Price as Given}\hfill{\large ECON 201}


\needspace{5\baselineskip}\section{Lena's lemonade}

Lena runs one of twenty lemonade stands at the county fair. All of them sell
the same cup of lemonade and have the same costs, and the going price is \$3 a
cup. The graph shows Lena's marginal cost: \$1 a cup for the first 60 cups, and
\$4 a cup after that, when she has to buy ice at the fair, up to 90 cups.

\begin{center}
\includegraphics[width=0.6\textwidth, alt={A plotting grid, price and cost on
the vertical axis from 0 to 6 dollars with gridlines every dollar, and cups of
lemonade on the horizontal axis from 0 to 100 with gridlines every 10 cups.
Marginal cost is a step: flat at 1 dollar from 0 to 60 cups, then a vertical
jump at 60 cups to 4 dollars, and flat at 4 dollars from 60 to 90
cups.}]{figures/lemonade.pdf}
\end{center}

\vspace{4pt}
1. Draw the demand curve that Lena faces on the graph. What happens if she
charges \$3.50?
\answerlines{1}

2. At \$3, how much does the 30th cup add to her profit? The 70th cup? How many
cups does she sell?
\answerlines{2}

3. Fill in the second column: how many cups does Lena sell at each going price?

\vspace{4pt}
4. All twenty stands have Lena's costs. Fill in the next column: how many cups
do the twenty stands sell in total at each price?

\begin{center}
\renewcommand{\arraystretch}{1.25}
\begin{tabular}{|c|C|C|C|}
\hline
\textbf{Going price, $P$} & \textbf{Cups Lena sells} & \textbf{Cups all twenty stands sell} & \textbf{Cups fairgoers want to buy} \\
\hline
\$0.50 & \blankrow & & 2,200 \\ \hline
\$2    & \blankrow & & 1,600 \\ \hline
\$3    & \blankrow & & 1,200 \\ \hline
\$5    & \blankrow & &   400 \\ \hline
\end{tabular}
\end{center}

5. Which curve on the graph tells you how many cups Lena sells at any going
price? What do we call it?
\answerlines{1}

6. The last column shows how many cups fairgoers want to buy at each price. At
which price does the market clear? How does it compare with the price Lena
takes as given?
\answerlines{2}

\end{document}
